\documentclass[11pt]{article}
\usepackage[a4paper,margin=27mm]{geometry}
\usepackage{amsmath,amssymb,mathtools}
\usepackage{fontspec,unicode-math}
\setmainfont{Latin Modern Roman}
\setmathfont{Latin Modern Math}
\usepackage{xurl,hyperref}
\hypersetup{colorlinks=true,urlcolor=blue,linkcolor=blue}
\newcommand{\sref}[2]{\hyperlink{source:#1:#2}{[S#2]}}
\newcommand{\cataloguescope}[1]{\paragraph{Mathematical question.} #1}
\begin{document}
\section{Multivariate analysis}\label{statement}

\subsection{Definitions and mathematical statement}

\paragraph{Shared notation.}$\mathbb N=\{1,2,\ldots\}$, and $\mathbb Z,\mathbb Q,\mathbb R,\mathbb C$ denote the integers, rationals, reals and complex numbers. Any other conventions are specified locally.


Let the $n$ columns $M_j\in\mathbb R^p$ be independent with law $N_p(\mu,\Sigma)$, where $\Sigma$ is positive definite with eigenvalues $\eta_1\ge\cdots\ge\eta_k>1$ and $\eta_{k+1}=\cdots=\eta_p=1$, with $k$ fixed as $p,n\to\infty$. Put $\bar M=n^{-1}\sum_j M_j$, $\widehat M_j=M_j-\bar M$, $S=n^{-1}\widehat M\widehat M^T$, and write $l_1\ge\cdots\ge l_p\ge0$ for the sample eigenvalues. Assume $p/n\to\gamma\in(0,\infty)$.
\paragraph{Statement recorded in the supplied source.}
Determine the asymptotic distributions of the leading sample eigenvalues $l_1,l_2,\ldots$ for these real spiked populations, identifying their appropriate centerings, scalings and joint laws throughout the different spike regimes, including the transition at $\eta=1+\sqrt\gamma$. \sref{statement}{2}
\subsection{Short English statement}
Find the limiting eigenvalue statistics for high-dimensional real Gaussian covariance matrices with finitely many population spikes, including phase transitions.
\subsection{Sources}
\begin{description}
\item[\hypertarget{source:statement:1}{S1}] ULAM.AI and the UnsolvedMath Contributors, UnsolvedMath: A Curated Collection of Open Mathematics Problems, catalogue entry. CC BY 4.0. \url{https://huggingface.co/datasets/ulamai/UnsolvedMath}.
\item[\hypertarget{source:statement:2}{S2}] Percy Deift, Some Open Problems in Random Matrix Theory and the Theory of Integrable Systems, arXiv:0712.0849 (2007), Problem 5, pp. 4--5. \url{https://arxiv.org/abs/0712.0849}.
\end{description}
\label{end:statement}
\end{document}
